%% ============================================================
%% Cross-Database Generalisation of ML-Driven Replication
%% Mode Selection — International Journal of Data Science
%% and Analytics (Springer)
%% Template: Springer Nature sn-jnl.cls (sn-mathphys, Numbered,
%%           iicol double-column layout per IJDSA guidelines)
%% Compiler : pdflatex
%% Status   : Final. All technical/reproducibility items resolved
%%            and verified against real data (see Table 3/5/6,
%%            Section 4.5). All author/advisor sign-offs confirmed.
%%            Ready for a final Overleaf compile and human read-through.
%% ============================================================

\documentclass[iicol,pdflatex,sn-mathphys-num]{sn-jnl}

% ---- Standard packages (per Journal of Grid Computing guidelines) ----
\usepackage{graphicx}
\usepackage{multirow}
\usepackage{amsmath,amssymb,amsfonts}
\usepackage{amsthm}
\usepackage{manyfoot}
\usepackage{booktabs}
\usepackage{xcolor}
\usepackage{textcomp}
\usepackage{url}
\usepackage{xurl}

% ---- Theorem-style environments (unused here, kept for consistency) ----
\theoremstyle{thmstyleone}
\newtheorem{theorem}{Theorem}
\newtheorem{proposition}[theorem]{Proposition}
\theoremstyle{thmstylethree}
\newtheorem{definition}{Definition}

\raggedbottom

%% ============================================================
\begin{document}

\title[Testing the Preconditions for Cross-Database Generalisation]{%
Testing the Preconditions for Cross-Database Generalisation of
Machine Learning-Driven Replication Mode Selection: A Controlled
Synthetic Study Across PostgreSQL, MySQL, and MongoDB}

\author*[1]{\fnm{Imane} \sur{Jerrari}}\email{imane.jerrari-etu@etu.univh2c.ma}

\author[1]{\fnm{Ismail} \sur{Assayad}}\email{ismail.assayad@univh2c.ma}

\affil*[1]{\orgdiv{Laboratory of Information Systems (LIS)},
\orgname{Faculty of Sciences Ain Chock \& ENSEM, Hassan II
University of Casablanca},
\orgaddress{\city{Casablanca}, \postcode{20000}, \country{Morocco}}}

\abstract{%
Our prior work introduced IAHA, a hybrid machine-learning framework
selecting among four replication modes (PERFORMANCE, CONSISTENCY,
AVAILABILITY, ECONOMIC) for containerised PostgreSQL, leaving
generalisation to other engines as future work. Here, under a
controlled design where all three offline datasets share a single
synthetic generator, we test whether a per-engine RandomForest
classifier behaves consistently across PostgreSQL, MySQL, and
MongoDB, and we validate live Kubernetes failover behaviour on all
three engines. Independent classifiers (PostgreSQL, $n=644$; MySQL,
$n=200$; MongoDB, $n=200$) significantly outperform a static
CPU-threshold heuristic by $+18.0$ to $+21.0$ percentage points
(McNemar $p<10^{-6}$), with mean accuracy $95.90\%$--$98.62\%$; a
direct comparison shows RandomForest matches
the strongest alternative (XGBoost) and both substantially outperform
linear/kernel classifiers. A size-matched control shows part of
PostgreSQL's advantage comes from dataset size (a $1.93$-point drop
at $n=200$), but size-matched PostgreSQL lands between, not below,
MySQL and MongoDB, so size alone does not fully explain cross-engine
differences.
Feature-importance rankings correlate strongly across engines
(Spearman $\rho=0.81$--$0.92$). All three classifiers misclassify
AVAILABILITY as CONSISTENCY; given small per-engine counts
($n=1,3,5$), we treat this descriptively; SMOTE improves recall only
to $20\%$. Live failover tests show total recovery time differs
significantly between engines (Mann--Whitney $p<0.05$), with
engine-specific differences and measurement artefacts reported
openly rather than resolved and confounded by an orchestration
asymmetry across engines. Because the
offline datasets share one generative process, those results
establish consistency under controlled conditions, not generalisation
to heterogeneous production systems; we discuss this distinction and
its implications for future work.}

\keywords{High Availability, Containerised Databases, Cross-Database
Generalisation, Machine Learning, Replication Mode Selection}

\maketitle

%% ============================================================
\section{Introduction}\label{sec:intro}
%% ============================================================

\subsection{Context and Motivation}\label{subsec:context}

% [Confirmed by author, 19 Jul 2026: iaha2026 status remains "under
%  review" at JIFS, matching \bibitem{iaha2026} as currently written.
%  No update needed unless the status changes before submission.]
Adaptive replication-mode selection has recently been shown to
outperform static heuristics on containerised PostgreSQL clusters
\cite{iaha2026}, using a RandomForest classifier trained on operational
metrics to choose among four modes at runtime. That study, however,
was scoped explicitly to a single database engine, and identified
transfer learning to MySQL and MongoDB as a direct item of future
work (FW4 in \cite{iaha2026}). IAHA operates in two layers: a fast,
rule-based P1 layer that applies static CPU-utilisation thresholds as
an immediate fallback, and a P2 layer in which a RandomForest
classifier is invoked at runtime on a ten-feature operational
snapshot (Table~\ref{tab:features}) to select among the four
replication modes. The P2 classifier is trained offline on
expert-validated operational traces and queried online at a fixed
polling interval, with mode transitions applied through
CloudNativePG's replication configuration; the static heuristic used
as this paper's own baseline (Section~\ref{subsec:stats}) mirrors
IAHA's P1 layer specifically, so that both this paper's comparison
and IAHA's own evaluation isolate the value the P2 classifier adds
over the simpler fallback it was designed to replace.

This leaves an open empirical question that matters beyond PostgreSQL:
is the relationship between operational metrics (CPU utilisation,
transaction throughput, replication lag, etc.) and the appropriate
replication mode a property of PostgreSQL specifically, or a more
general property of replicated database systems under load? If the
latter, a single engine-agnostic feature schema could support adaptive
high availability (HA) across heterogeneous database fleets --- a
common scenario in practice, where organisations operate relational
and NoSQL engines side by side. As a first step toward answering that
question, this paper asks a narrower, more tractable one: under
controlled conditions that hold the data-generating process fixed
across engines, does classifier behaviour at least remain consistent?
A negative answer would undercut the premise of cross-engine transfer
before any live multi-engine validation is attempted; a positive
answer, while not proof of generalisation, would justify pursuing it.

\subsection{Problem Statement}\label{subsec:problem}

Under a controlled synthetic design that holds the data-generating
process fixed across engines, does a machine-learning classifier for
replication mode selection, trained independently on PostgreSQL,
MySQL, and MongoDB using an identical engine-agnostic feature schema,
achieve consistent classification performance across engines, and do
the same classes remain difficult to predict regardless of the
underlying engine?

\subsection{Research Questions}\label{subsec:rqs}

\newcommand{\rqlabel}[1]{\textbf{#1}\hspace{0.6em}---\hspace{0.4em}}

\begin{itemize}
\item[] \rqlabel{RQ1} Under a controlled synthetic design sharing a
  single generative process, is classifier performance (accuracy,
  macro-F1) consistent across PostgreSQL, MySQL, and MongoDB, or does
  it depend on engine-specific factors such as dataset size and class
  distribution?
\item[] \rqlabel{RQ1c} Is RandomForest, inherited from \cite{iaha2026},
  actually competitive with alternative classifiers on this specific
  four-class mode-selection task, or would another model perform
  better?
\item[] \rqlabel{RQ2} Is the relative importance of operational features
  for mode selection stable across a relational engine (PostgreSQL),
  a second relational engine (MySQL), and a document-oriented engine
  (MongoDB)?
\item[] \rqlabel{RQ3} Does the AVAILABILITY-class failure mode identified
  as a limitation in \cite{iaha2026} (L3: extreme class imbalance)
  show the same error pattern across engines, and is it consistent
  with a correctable sample-size artefact or with a structural
  feature-space overlap with CONSISTENCY? (Given the very small
  per-engine AVAILABILITY counts, we treat this as an exploratory,
  hypothesis-generating question rather than one this study is
  powered to settle definitively.)
\item[] \rqlabel{RQ4} Beyond offline classification accuracy, do live
  Kubernetes failover tests show consistent end-to-end recovery
  behaviour across the three engines, and what confounds (deployment
  orchestration, measurement artefacts, temporal effects) must be
  accounted for in interpreting any differences observed?
\end{itemize}

\subsection{Contributions}\label{subsec:contributions}

\begin{itemize}
\item A controlled synthetic evaluation of replication-mode
  classification across three database engines (PostgreSQL, MySQL,
  MongoDB), generated by a single shared engine-agnostic generator and
  feature schema, isolating whether classification behaviour depends
  on engine identity. We are explicit that this design tests a
  precondition for cross-engine generalisation --- consistency of
  classifier behaviour under a shared generative process --- rather
  than generalisation to heterogeneous production systems itself
  (Section~\ref{subsec:scope}).
\item A sample-size control experiment that directly disentangles
  engine identity from training-set size, showing that the apparent
  performance advantage of the largest dataset (PostgreSQL) is
  \emph{partly} --- but not fully --- attributable to its size:
  size-matched PostgreSQL lands between MySQL and MongoDB rather than
  below both, so a residual, unattributed difference between engines
  remains and is flagged as an open question rather than resolved.
\item Statistically grounded evidence (BCa bootstrap confidence
  intervals, McNemar tests against a static heuristic baseline) that
  ML-driven mode selection significantly outperforms static thresholds
  on all three engines, not only PostgreSQL: a $+18.0$ to $+21.0$
  percentage-point accuracy improvement, significant at
  $p<10^{-6}$ on every engine (Section~\ref{subsec:rq1}). This is the
  most unambiguous result in the paper and holds regardless of how
  the other, more exploratory findings below are read.
\item A direct comparison of RandomForest against XGBoost, logistic
  regression, and RBF-SVM on the same task (Section~\ref{subsec:rq1c}),
  showing RandomForest is empirically competitive with the strongest
  alternative (XGBoost) rather than an unexamined default, and that
  both tree-based methods substantially outperform linear/kernel
  methods --- a result requiring no hedging.
\item Evidence of cross-engine stability in feature importance rankings
  (Spearman $\rho \geq 0.81$), motivating future transfer-learning work.
\item A live Kubernetes failover validation across all three engines
  (Section~\ref{subsec:rq4}), extending the offline classification
  results with measured end-to-end recovery times. We report this
  alongside, rather than instead of, the limitations we uncovered in
  the process: an orchestration asymmetry between engines, a
  temporal ordering effect in the PostgreSQL tests, and an unresolved
  measurement artefact in a subset of the MongoDB tests --- so that
  the live results are read as an initial validation rather than a
  fully controlled comparison.
\end{itemize}

%% ============================================================
\section{Related Work}\label{sec:related}
%% ============================================================

\subsection{Adaptive High Availability for Databases}\label{subsec:ha}

Existing PostgreSQL HA tools follow a common pattern: a single
synchronous or asynchronous replication policy chosen once at
deployment time and left unchanged at runtime. Patroni \cite{patroni}
and repmgr \cite{repmgr} provide consensus-based leader election but
do not adapt replication behaviour to load; Stolon \cite{stolon}
follows the same fixed-policy pattern. CloudNativePG \cite{cloudnativepg},
the CNCF-donated Kubernetes operator used as the substrate
in \cite{iaha2026}, improves failover latency through a native
Kubernetes control loop but likewise applies one static replication
configuration for the cluster's entire lifetime. None of these tools
offers more than two operational modes, and none adapts mode selection
to observed workload conditions --- the gap that IAHA \cite{iaha2026}
addresses for PostgreSQL, and that this paper extends to MySQL and
MongoDB.

\subsection{Machine Learning for Database Management}\label{subsec:ml4db}

ML-driven database optimisation is an active area spanning several
problem classes: cardinality estimation and query optimisation
(SageDB \cite{kraska2019}, Neo \cite{marcus2019}, Bao \cite{marcus2021}),
automated configuration tuning
(OtterTune \cite{vanaken2017}, CDBTune \cite{zhang2019}), and resource
management (Autopilot \cite{autopilot2020}, FIRM \cite{firm2020}). A
recurring finding in this literature, most directly relevant to our
RQ1, is that ML-driven tuning gains do \emph{not} automatically
transfer across systems or even across versions of the same system:
OtterTune \cite{vanaken2017} explicitly maps unseen workloads to
previously characterised ones because tuning effort on one DBMS does
not reduce tuning effort on another, and LlamaTune \cite{kanellis2022}
reports that the throughput improvement of its tuning approach over a
baseline optimiser narrows substantially when ported from PostgreSQL
v9.6 to v13.6 on certain workloads (e.g.\ from $20.8\%$ to $3.6\%$ on
YCSB-B), despite both being the same engine. Multi-DBMS evaluations
such as GPTuner \cite{lao2024} and UDO \cite{wang2021udo}, which
benchmark across PostgreSQL and MySQL, similarly report
engine-dependent variation in tuning effectiveness rather than uniform
gains. This body of work motivates our explicit empirical test (RQ1,
RQ1b) of whether classification performance is consistent across
engines or is confounded by other factors --- in our case, training-set
size --- rather than assuming consistency holds by default. None of
these systems addresses replication-mode selection specifically, nor
provides operator-facing explainability, which remains the gap
IAHA \cite{iaha2026} and this paper target.

\subsection{Explainable AI in Critical Systems}\label{subsec:xai}

Explainable AI (XAI) addresses the interpretability of ML-driven
decisions \cite{gunning2019} through post-hoc methods such as
SHAP \cite{lundberg2017} and LIME \cite{ribeiro2016}, counterfactual
explanations \cite{wachter2017}, and rule extraction \cite{guidotti2018}.
XAI has been applied to high-stakes domains including
healthcare \cite{holzinger2019}, autonomous vehicles \cite{atakishiyev2021},
and finance \cite{bussmann2020}, but --- as noted in \cite{iaha2026}
--- its application to database HA management remains largely
unexplored. This paper does not introduce new XAI components beyond
those already proposed in \cite{iaha2026}; instead, the feature
importance analysis in Section~\ref{subsec:rq2} serves a
complementary diagnostic role, testing whether the same decision
signature used for explainability on PostgreSQL is structurally stable
enough to remain meaningful on MySQL and MongoDB.

\subsection{Cross-Database and Transfer Learning Approaches}\label{subsec:crossdb}

Despite extensive work on ML for single-DBMS optimisation
(Section~\ref{subsec:ml4db}), explicit cross-engine generalisation
remains comparatively underexplored. Most multi-DBMS evaluations in
the tuning literature (e.g.\ GPTuner \cite{lao2024},
UDO \cite{wang2021udo}) train and evaluate \emph{separate} models per
engine to demonstrate applicability, rather than testing whether a
model or feature representation \emph{transfers} from one engine to
another --- the same design choice we make in this paper (independent
per-engine classifiers, Section~\ref{sec:methodology}), motivated by
the same practical consideration: heterogeneous internal metrics and
configuration spaces across engines make naive model reuse
non-trivial. Industrial efforts to copy or retrain ML models across
database instances exist as systems-engineering patterns, but
typically address operational mechanics of model transfer (e.g.\
avoiding redundant retraining) rather than evaluating whether the
\emph{learned relationships} generalise across heterogeneous engines.
To our knowledge, no prior work evaluates whether the operational
signals driving a replication-mode mapping learned on one database
engine are structurally consistent with those of a different engine
family (relational to relational, or relational to document-oriented)
--- the specific gap this paper addresses empirically through RQ1--RQ3.
We note that this paper tests the \emph{preconditions} for transfer
(consistent classification behaviour, stable feature-importance
signatures) rather than transfer itself; directly training on one
engine and testing on another is identified as future work (FW1).

\textbf{Why the systems cited above are not empirical baselines
here.} Section~\ref{subsec:ml4db} discusses OtterTune, CDBTune,
GPTuner, UDO, and LlamaTune at length because they are the closest
prior work on ML-driven database management, and we return to them
here to state explicitly why none is used as an empirical comparison
point in this paper's results. These systems perform configuration
\emph{tuning} --- selecting numerical parameter values (buffer sizes,
timeouts, memory allocations) to optimise a continuous performance
objective --- whereas this paper performs replication-mode
\emph{selection}, a four-class classification problem over discrete
operational states. Patroni, repmgr, and CloudNativePG
(Section~\ref{subsec:ha}) are HA orchestration tools that execute a
fixed failover policy rather than choosing among multiple named
replication modes at all. Running any of these systems on our task
would not produce a comparable output: a tuning system would need to
be given a four-class decision problem it was not designed to solve,
and an orchestration tool would need a mode-selection layer bolted on
that does not exist in its original design. An empirical comparison
under these conditions would not isolate a meaningful difference in
method --- it would compare our classifier against our own
re-implementation of a decision layer these systems do not have,
which we do not consider informative. We therefore restrict empirical
comparison in this paper to (i) the static heuristic baseline
(Section~\ref{subsec:stats}), which performs the same classification
task under a simpler decision rule, and (ii) alternative ML
classifiers on the same task (Section~\ref{subsec:rq1c}), rather than
attempting a system-level comparison that would not be apples-to-apples.

\subsection{Research Gaps Addressed by This Paper}\label{subsec:gaps}

Building on the four gaps identified in \cite{iaha2026} (binary modes,
static thresholds, no explainability, no temporal optimisation), we
identify a fifth gap specific to cross-engine generalisation:

\textbf{Gap 5 -- Untested Cross-Engine Generalisation:} Existing
ML-driven database management systems are evaluated on a single engine
(or, where multiple engines are used, evaluated independently without
testing whether learned relationships transfer between them), leaving
open whether reported performance gains reflect properties of the
engine, the dataset, or the underlying operational-metric relationships
themselves. This paper addresses Gap 5 through the size-matched
control analysis in Section~\ref{subsec:rq1b}, which shows dataset
size is part, but not all, of the explanation for the cross-engine
differences observed: it does not fully close the gap, since
size-matched PostgreSQL lands between MySQL and MongoDB rather than
below both, leaving a residual difference between engines at matched
sample size open (Section~\ref{subsec:rq1b}, FW4), and --- as
discussed in Section~\ref{subsec:scope} --- closing it fully will
additionally require moving beyond a single shared synthetic generator
to independently instrumented engines.

%% ============================================================
\section{Methodology}\label{sec:methodology}
%% ============================================================

\subsection{Datasets}\label{subsec:datasets}

Three datasets were used, sharing an identical ten-feature schema
(Table~\ref{tab:features}) and four-class label space (PERFORMANCE,
CONSISTENCY, AVAILABILITY, ECONOMIC). All three are produced by a
\emph{single shared synthetic generator} (a common
\texttt{generate\_samples(n, system)} Python function, confirmed by
direct inspection of the generation script), called three times with
the same random seed and the same scenario-based labelling logic,
varying only the requested sample size $n$ and an engine-specific
parameter \texttt{system}. Concretely, each sample is drawn from one
of five latent operating scenarios (normal, peak, off-peak, crisis,
degraded), and the four-class mode label is then assigned by a
deterministic rule cascade over the sampled feature values (load
thresholds, error rates, deadlock counts, and time-of-day); this is
what produces the class distributions in Table~\ref{tab:class-dist}
below, and the full mechanical detail (scenario probabilities,
parameter ranges, and the exact rule cascade) is given in Limitation
L1 (Section~\ref{subsec:limitations}) rather than repeated here. We
state this design choice here, rather than only in Limitation L1,
because it directly shapes how the
cross-engine comparisons in Section~\ref{sec:results} should be read:
the three datasets are controlled parametric variations of one
generative process, not independent instrumentation of three
production systems. This guarantees a consistent, comparable feature
schema by construction, but it also means the study should be
understood as a controlled synthetic evaluation, not an empirical
survey of real production traces (see Limitation L1 for the full
generator detail and Threats to Validity for the consequences).

\subsection{Scope: What This Study Does and Does Not Show}\label{subsec:scope}

Because this design choice is central to interpreting every result
that follows, we state it plainly before presenting the feature
schema and class distributions. This study shows that a RandomForest
classifier, trained independently per engine, behaves
\emph{consistently} --- in accuracy, in error pattern, and in which
features it relies on --- across three parametric variations of a
single synthetic generator that differ only in requested sample size
and a small number of engine-specific parameter ranges
(Section~\ref{subsec:datasets}; Limitation L1). This consistency is a
\emph{necessary precondition} for cross-database transfer learning to
be worth pursuing: if a shared generative process produced
inconsistent classifier behaviour across engines, there would be
little reason to expect a model trained on one real engine to say
anything useful about another.

This study does \emph{not} show that a classifier trained on
PostgreSQL data generalises to independently instrumented, live
MySQL or MongoDB deployments. The three datasets do not capture
engine-specific implementation details, workload idiosyncrasies, or
telemetry artefacts that would appear in production traces collected
separately from each engine; two features
(\texttt{wal\_buffers\_full}, \texttt{trans\_rolled\_back}) vary in
scale between engines within the generator, but the remaining eight,
and the entire labelling-rule cascade, are identical by construction.
Consequently, results such as the high cross-engine Spearman
correlation in feature importance (Section~\ref{subsec:rq2}) should
be read as evidence that the \emph{signature} of a good classifier is
stable under this controlled design, not as evidence that the
signature would remain stable if measured on independently
instrumented production systems. We treat genuine cross-engine
transfer --- training on one real engine and evaluating on
another --- as the open question this paper motivates but does not
answer, and identify it explicitly as future work (FW1,
Section~\ref{sec:conclusion}).

\begin{itemize}
\item \textbf{PostgreSQL}: $n = 644$ expert-validated samples
  (the same dataset used to train IAHA's Phase 2 classifier
  in \cite{iaha2026}).
\item \textbf{MySQL}: $n = 200$ samples, generated under the same
  schema for this study.
\item \textbf{MongoDB}: $n = 200$ samples, generated under the same
  schema for this study.
\end{itemize}

We did not generate $n=644$ samples for MySQL and MongoDB to match
PostgreSQL, and we state the reason explicitly rather than leaving it
implicit: the PostgreSQL dataset is not newly generated for this
paper but is reused, unchanged, from \cite{iaha2026}, where it
underwent independent development and informal validation over the
course of that earlier study. Generating a matching $n=644$ sample
for MySQL and MongoDB would not reproduce that same provenance --- it
would simply be more synthetic data from the same generator, which
would not address the underlying asymmetry (real prior validation
history vs.\ newly generated data) and could give a misleading
impression of equivalence between the three datasets. We instead
treat the size difference as a variable to be explicitly controlled
for, which is exactly the role of the size-matched analysis in
Section~\ref{subsec:rq1b}, rather than a difference to be papered
over by generating more same-generator data for the other two
engines.

\begin{table}[h]
\caption{Shared Feature Schema Across Engines}\label{tab:features}
\begin{tabular}{@{}ll@{}}
\toprule
Feature & Description \\
\midrule
\texttt{replication\_lag\_mb} & Replication lag (MB) \\
\texttt{transactions\_per\_sec} & Transaction throughput \\
\texttt{cpu\_percent} & CPU utilisation (\%) \\
\texttt{latency\_ms} & Query/operation latency \\
\texttt{trans\_rolled\_back} & Rolled-back transaction rate \\
\texttt{deadlocks} & Deadlock count \\
\texttt{hour\_of\_day} & Time-of-day indicator \\
\texttt{memory\_usage\_pct} & Memory utilisation (\%) \\
\texttt{connections\_active} & Active connection count \\
\texttt{wal\_buffers\_full} & Write-ahead/journal buffer saturation \\
\botrule
\end{tabular}
\end{table}

\textbf{Engine-agnostic feature design.} Features were deliberately
chosen as generic operational metrics rather than PostgreSQL-specific
internals, to enable this cross-engine comparison without
re-engineering the feature set per database. We note that
\texttt{wal\_buffers\_full} is measured as a generic write-ahead/journal
buffer saturation proxy; for MongoDB this corresponds functionally to
oplog buffer pressure rather than a literal PostgreSQL WAL buffer, and
we report it under this generic name rather than renaming it
per-engine, to preserve direct cross-engine comparability of the same
underlying measurement axis. Similarly, \texttt{deadlocks} and
\texttt{trans\_rolled\_back} are reported as generic
concurrency-conflict and transaction-failure indicators rather than
engine-native terminology.

\textbf{Synthetic data and class imbalance.} As in \cite{iaha2026}
(Limitation L2), all three datasets are synthetically generated under
expert-informed thresholds rather than collected from production
traces. Class distributions are imbalanced across all three engines,
with CONSISTENCY accounting for $85.0\%$ (MongoDB), $78.5\%$ (MySQL),
and $79.5\%$ (PostgreSQL) of samples, and AVAILABILITY represented by
only $n=1$, $n=3$, and $n=5$ samples respectively
(Table~\ref{tab:class-dist}).

\begin{table}[h]
\caption{Class Distribution by Engine}\label{tab:class-dist}
\begin{tabular}{@{}lrrr@{}}
\toprule
Mode & MongoDB & MySQL & PostgreSQL \\
\midrule
CONSISTENCY  & 170 (85.0\%) & 157 (78.5\%) & 512 (79.5\%) \\
ECONOMIC     & 17 (8.5\%)   & 22 (11.0\%)  & 57 (8.9\%)   \\
PERFORMANCE  & 12 (6.0\%)   & 18 (9.0\%)   & 70 (10.9\%)  \\
AVAILABILITY & 1 (0.5\%)    & 3 (1.5\%)    & 5 (0.8\%)    \\
\midrule
Total ($n$)  & 200 & 200 & 644 \\
\botrule
\end{tabular}
\end{table}

\subsection{Classification Protocol}\label{subsec:protocol}

For each engine, we train an independent RandomForest classifier
(200 trees, balanced class weighting, \texttt{random\_state=42} for
both the classifier and the cross-validation fold assignment, fixed
throughout all experiments in this paper for reproducibility) using
repeated stratified 5-fold cross-validation (10 repeats, 50 fold-level
accuracy estimates per engine), matching the cross-validation protocol
used in \cite{iaha2026}.

\subsection{Alternative Classifier Comparison}\label{subsec:classifier-comparison-protocol}

As justified at length in Section~\ref{subsec:crossdb}, we
restrict empirical comparison in this paper to the static heuristic
baseline (Section~\ref{subsec:stats}) and to alternative ML
classifiers on the same task (this section), rather than to the
configuration-tuning systems or HA orchestration tools discussed in
Section~\ref{subsec:ml4db} and Section~\ref{subsec:ha},
since none of those systems can be run on this four-class
mode-selection task without a redesign that would make any comparison
uninformative. RandomForest is used throughout this paper because it
was the model selected in \cite{iaha2026}, not because it was
benchmarked against
alternatives for this specific four-class mode-selection task. To
address this directly, we additionally train three alternative
classifiers under the identical protocol (same features, same
repeated $10\times5$-fold stratified CV, same \texttt{random\_state=42}):
XGBoost (200 trees; all other hyperparameters left at
xgboost~3.4.1's library defaults, i.e.\ no additional tuning beyond
matching RandomForest's tree count), a logistic regression
(\texttt{class\_weight="balanced"}, \texttt{max\_iter=1000}, default
\texttt{lbfgs} solver), and an RBF-kernel SVM
(\texttt{class\_weight="balanced"}, default kernel coefficient), both
with balanced class weighting where supported
and with features standardised (zero mean, unit variance) for the
two distance/gradient-based models. All four classifiers use
\texttt{random\_state=42}, matching the seed used everywhere else in
this paper. This does not compare our
classifier against the ML-driven tuning systems discussed in
Section~\ref{subsec:ml4db} (OtterTune, GPTuner, and similar), since,
as discussed in Section~\ref{subsec:crossdb}, those systems solve a
different task (continuous parameter tuning) and cannot be run on a
four-class mode-selection problem without redesign; it instead
answers the narrower and directly answerable question of whether
RandomForest is competitive with other standard classifiers on
exactly this task.

\subsection{Statistical Testing}\label{subsec:stats}

\textbf{BCa bootstrap confidence intervals.} Following the
methodology of \cite{iaha2026}, we report bias-corrected and
accelerated (BCa) bootstrap 95\% confidence intervals (10{,}000
resamples) on mean cross-validated accuracy per engine.

\textbf{McNemar test.} To test whether the RandomForest classifier
significantly outperforms a static heuristic, we apply McNemar's test
\emph{within each engine}, comparing paired correct/incorrect
predictions on the same held-out folds. The static heuristic mirrors
the threshold logic of IAHA's P1 layer: \texttt{cpu\_percent}\,$<30 \Rightarrow$
ECONOMIC; $30\leq$\texttt{cpu\_percent}\,$<50 \Rightarrow$ PERFORMANCE;
otherwise CONSISTENCY (the heuristic cannot express AVAILABILITY).
We do \emph{not} apply McNemar across engines (e.g.\ MongoDB vs.\
MySQL), since the two samples are drawn from different observation
spaces and are not paired --- this would violate the test's
matched-pairs assumption.

\textbf{Mann--Whitney U test.} To compare the distributions of
cross-validated accuracy between pairs of engines, we use the
Mann--Whitney U test on the 50 fold-level accuracy estimates per
engine. We note as a limitation that these 50 estimates per engine are
not fully independent (they derive from repeated resampling of the
same underlying dataset), so the resulting $p$-values likely overstate
the true level of statistical confidence; we report them as indicative
rather than definitive evidence of inter-engine difference.

\subsection{Sample-Size Control Protocol}\label{subsec:sizecontrol-protocol}

Because PostgreSQL ($n=644$) is $3.2\times$ larger than MySQL or
MongoDB ($n=200$ each), any raw accuracy comparison between engines is
confounded with a sample-size effect. To isolate the contribution of
engine identity from dataset size, we additionally evaluate a
\emph{size-matched} PostgreSQL condition: 20 independent stratified
subsamples of $n=200$ are drawn from the full PostgreSQL dataset
(class-proportional: 158 CONSISTENCY, 22 PERFORMANCE, 18 ECONOMIC, 2
AVAILABILITY), each evaluated under the same repeated stratified
cross-validation protocol as the full dataset (10 repeats $\times$ 5
folds), yielding $20 \times 50 = 1000$ fold-level accuracy estimates
in total. This size-matched PostgreSQL condition is then
compared against (i) PostgreSQL at full size, and (ii) MySQL and
MongoDB at their native $n=200$, using the same Mann--Whitney
procedure described above.

\subsection{Live Multi-Engine Cluster Validation}\label{subsec:live-validation}

In addition to the offline classification evaluation above, we
conducted live failover testing on containerised clusters for all
three engines, to measure end-to-end recovery time rather than only
classification accuracy. This directly extends the Phase~1
CloudNativePG validation of \cite{iaha2026} to MySQL and MongoDB.

\textbf{Cluster setup and an explicit orchestration asymmetry.}
PostgreSQL was deployed via CloudNativePG v1.22.0 (PostgreSQL 16.1,
3 instances), the same production-grade Kubernetes operator used
in \cite{iaha2026}. MySQL was deployed via the MySQL Operator's
\texttt{InnoDBCluster} resource (3 instances, Group Replication, 1
router instance). MongoDB was deployed as a bare 3-replica
\texttt{StatefulSet} running \texttt{mongod --replSet} directly
(MongoDB Community Server 7.0.0), \emph{without} a dedicated MongoDB
Kubernetes operator. We state this asymmetry explicitly rather than
letting it pass unremarked: PostgreSQL and MySQL benefit from
operator-managed orchestration (automated leader election tooling,
health probes, and reconciliation loops built on top of each engine's
native replication), while MongoDB's failover in this setup relies
solely on the database's own internal replica-set election with no
additional Kubernetes-level orchestration. This is a threat to the
fairness of the cross-engine RTO comparison (Section~\ref{subsec:rq4})
and is discussed further in Threats to Validity
(Section~\ref{subsec:threats}); it means the comparison below is
better read as ``CloudNativePG-managed PostgreSQL vs.\
Operator-managed MySQL vs.\ unmanaged MongoDB'' rather than as a
clean, orchestration-neutral comparison of the three database
engines themselves.

\textbf{Experimental environment.} All live cluster tests and offline
classification experiments were run on a single host under Windows
11 with WSL2 (Ubuntu), using Docker Desktop with the WSL2 backend and
\texttt{kind} (Kubernetes-in-Docker) to provision one local cluster
per engine (\texttt{kind-iaha-postgres}, \texttt{kind-iaha-mysql},
\texttt{kind-iaha-mongodb}). Host specification: Intel(R) Core(TM)
i5-8250U CPU @ 1.60GHz (4 physical cores, 8 logical processors),
16\,GB RAM; Docker
Desktop version 28.3.2 (build 578ccf6); \texttt{kind} version 0.20.0.
No other CPU-intensive
process was intentionally run on the host during data collection,
though WSL2's shared-kernel virtualisation means host-level
scheduling noise cannot be fully ruled out as a contributor to the
timing variability discussed in Section~\ref{subsec:rq4} and
Section~\ref{subsec:threats}. The offline classification experiments
(Sections~\ref{subsec:protocol}--\ref{subsec:stats}) were run under
Python 3.13.14 with the package versions pinned in
\texttt{requirements-lock.txt} in the accompanying repository, rather
than the minimum versions given in \texttt{requirements.txt}, so that
the published numbers can be reproduced exactly rather than only
approximately.

\textbf{Failover protocol.} For each test, we identify the current
primary/leader pod, delete it via \texttt{kubectl delete pod}, and
measure two sequential, non-overlapping intervals: (i)
\texttt{promotion\_time\_s}, the time from pod deletion until a new
primary is detected, and (ii) a post-promotion stabilisation interval
until the cluster reports all replicas healthy/online again. We refer
to the sum of these two intervals as \emph{total RTO}
($\text{total\_rto\_s} = \text{promotion\_time\_s} + \text{stabilisation\_s}$)
throughout this paper. We flag this definition explicitly because our
measurement scripts label the second interval alone as
\texttt{rto\_s}; read in isolation, that field understates true
end-to-end recovery time, since it excludes the promotion-detection
phase that necessarily precedes it. All RTO figures reported below
use the summed, total quantity.

\textbf{Sample sizes and exclusions.} PostgreSQL: 30/30 attempted
failover tests completed successfully. MySQL: 15 tests were attempted
before testing was halted after test 15 failed to reach 3 ONLINE
Group Replication members within the 180s timeout (recorded as a
sentinel value in the raw log); this left $n=14$ valid measurements.
MongoDB: 30/30 tests returned a result, but 10 of the 30 produced a
stabilisation time tightly clustered around 0.65--0.73s -- implausibly
fast and implausibly uniform compared to the remaining 20 tests
(4.36--19.39s) -- which we judge to most likely reflect a measurement
artefact (e.g.\ a stale replica-set status read immediately after pod
deletion, before Kubernetes had actually evicted the terminating pod)
rather than genuine sub-second recovery; we therefore report MongoDB
results on the remaining $n=20$ tests and flag the excluded 10 as an
unresolved measurement question rather than silently discarding them
without explanation. We were unable to conclusively confirm the root
cause of this artefact and note it as a limitation
(Section~\ref{subsec:limitations}).

%% ============================================================
\section{Results}\label{sec:results}
%% ============================================================

\subsection{RQ1: Cross-Engine Classification Performance}\label{subsec:rq1}

Table~\ref{tab:main-results} summarises classification performance
across the three engines.

\begin{table}[h]
\caption{Classification Performance by Engine (Repeated $10\times5$-Fold CV)}\label{tab:main-results}
\begin{tabular}{@{}lrrr@{}}
\toprule
Metric & MongoDB & MySQL & PostgreSQL \\
\midrule
Mean accuracy & 97.55\% & 95.90\% & 98.62\% \\
BCa 95\% CI & [96.9, 98.1] & [95.3, 96.6] & [98.4, 98.8] \\
Majority baseline & 85.00\% & 78.50\% & 79.50\% \\
Heuristic acc. & 79.00\% & 75.00\% & 79.19\% \\
Macro-F1 (3 cl.)\footnotemark[1] & 93.15\% & 93.83\% & 98.93\% \\
\botrule
\end{tabular}
\footnotetext[1]{CONSISTENCY, ECONOMIC, PERFORMANCE only; AVAILABILITY
excluded from aggregate macro-F1 due to insufficient sample size for
reliable estimation (see Section~\ref{subsec:rq3}). Computed by
pooling predictions across all 50 repeated-CV folds (10 repeats
$\times$ 5 folds) and taking a single macro-averaged F1 over the
three classes on the pooled predictions, using scikit-learn's
\texttt{f1\_score(..., average="macro", zero\_division=0)}; this
matches \texttt{src/train\_classifiers.py} in the accompanying
repository exactly.}
\end{table}

All three classifiers substantially exceed both the majority-class
baseline and the static heuristic. We note that for PostgreSQL the
static heuristic ($79.19\%$) slightly underperforms the majority-class
baseline ($79.50\%$): the heuristic actively misclassifies some
CONSISTENCY samples into ECONOMIC or PERFORMANCE when
\texttt{cpu\_percent} crosses its thresholds, whereas a baseline that
always predicts CONSISTENCY incurs no such cost; this is consistent
with the heuristic providing no information about AVAILABILITY
and a CPU threshold that only weakly separates CONSISTENCY from
PERFORMANCE (Section~\ref{subsec:rq3}). The RandomForest classifier
does not exhibit this issue on any engine. McNemar's test confirms the
RandomForest classifier significantly outperforms the static
heuristic on every engine (Table~\ref{tab:mcnemar}). Note that the RF
accuracy reported in Table~\ref{tab:mcnemar} is computed on a single
stratified 5-fold split (required for McNemar's paired-prediction
design) rather than the repeated $10\times5$-fold protocol used for
Table~\ref{tab:main-results}; the two are not directly comparable and
the small numerical difference (e.g.\ $98.9\%$ vs.\ $98.62\%$ for
PostgreSQL) reflects this difference in protocol, not an
inconsistency in the underlying model.

\begin{table}[h]
\caption{McNemar Test: RandomForest vs.\ Static Heuristic (Per Engine)}\label{tab:mcnemar}
\begin{tabular}{@{}lrrrr@{}}
\toprule
Engine & RF acc. & Heur.\ acc. & $\Delta$ (pp) & McNemar $p$ \\
\midrule
MongoDB     & 97.0\% & 79.0\% & $+18.0$ & $2.5\times10^{-7}$ \\
MySQL       & 96.0\% & 75.0\% & $+21.0$ & $1.3\times10^{-8}$ \\
PostgreSQL  & 98.9\% & 79.2\% & $+19.7$ & $3.5\times10^{-28}$ \\
\botrule
\end{tabular}
\end{table}

The magnitude of improvement ($+18.0$ to $+21.0$ percentage points) is
consistent with the $+25.7$ percentage-point improvement over static
heuristics reported for PostgreSQL in \cite{iaha2026}, independently
confirming on two additional engines that ML-driven mode selection
provides a substantial, statistically significant advantage over fixed
thresholds.

\begin{figure*}[t]
\centering
\includegraphics[width=0.85\linewidth]{figures/fig1_accuracy_vs_heuristic.pdf}
\caption{RandomForest classifier accuracy versus the static
CPU-threshold heuristic, by engine. Error bars show BCa 95\%
confidence intervals on the repeated $10\times5$-fold CV accuracy
(Table~\ref{tab:main-results}); heuristic bars correspond to
Table~\ref{tab:mcnemar}.}\label{fig:accuracy-heuristic}
\end{figure*}

\textbf{Per-class precision and recall.} Table~\ref{tab:main-results}
and Table~\ref{tab:mcnemar} report only aggregate accuracy and
macro-F1; Table~\ref{tab:per-class} completes the picture with
per-class precision and recall for all four classes, computed on the
same pooled predictions (all 50 repeated-CV folds concatenated) used
for the pooled macro-F1 in Table~\ref{tab:main-results}, so the two
are directly consistent. CONSISTENCY and ECONOMIC are both recovered
with precision and recall at or above $95\%$ on every engine
(ECONOMIC reaching $100\%$ throughout, consistent with its empirical
separability discussed in Section~\ref{subsec:economic-note});
PERFORMANCE recall is lower and more variable across engines
($95.3\%$ PostgreSQL, $71.1\%$ MySQL, $67.5\%$ MongoDB), consistent
with the training-set-size effect on this class discussed in
Section~\ref{subsec:rq1b}; AVAILABILITY recall is $0\%$ on every
engine, discussed separately in Section~\ref{subsec:rq3} given the
very small per-engine sample sizes involved.

\begin{table*}[t]
\caption{Per-Class Precision and Recall (Pooled Predictions, All 50
Repeated-CV Folds)}\label{tab:per-class}
\begin{tabular}{@{}lrrrrrr@{}}
\toprule
 & \multicolumn{2}{c}{MongoDB} & \multicolumn{2}{c}{MySQL} & \multicolumn{2}{c}{PostgreSQL} \\
Class & Prec.\ & Rec.\ & Prec.\ & Rec.\ & Prec.\ & Rec.\ \\
\midrule
CONSISTENCY  & 97.2\%  & 100.0\% & 95.0\%  & 100.0\% & 98.4\%  & 99.9\%  \\
ECONOMIC     & 100.0\% & 100.0\% & 100.0\% & 100.0\% & 100.0\% & 100.0\% \\
PERFORMANCE  & 100.0\% & 67.5\%  & 100.0\% & 71.1\%  & 99.1\%  & 95.3\%  \\
AVAILABILITY & 0.0\%   & 0.0\%   & 0.0\%   & 0.0\%   & 0.0\%   & 0.0\%   \\
\botrule
\end{tabular}
\end{table*}

\textbf{Engine-level differences.} A Mann--Whitney U test on
fold-level accuracy distributions found all three pairwise comparisons
statistically significant (MongoDB vs.\ MySQL: $U=1739.0$,
$p=3.6\times10^{-4}$; MongoDB vs.\ PostgreSQL: $U=798.0$, $p=0.0016$;
MySQL vs.\ PostgreSQL: $U=240.0$, $p=1.9\times10^{-12}$), though as
discussed in
Section~\ref{subsec:stats} these $p$-values should be read as
indicative rather than definitive given the non-independence of
repeated-CV estimates. We do not apply a formal multiple-comparison
correction across the pairwise tests reported in this section and in
Sections~\ref{subsec:rq1b} and~\ref{subsec:rq4}; however, even a
conservative Bonferroni correction across all pairwise comparisons
reported in this paper would not change any conclusion, since the
largest $p$-value among them is $0.047$ (Section~\ref{subsec:rq4}).
PostgreSQL --- the largest dataset
($n=644$) --- achieves the highest accuracy and the tightest
confidence interval, which raises an immediate confound already
flagged in Section~\ref{sec:methodology}: is PostgreSQL's apparent
advantage a property of the engine, or simply a property of having
$3.2\times$ more training data? We test this directly in
Section~\ref{subsec:rq1b}.

\subsection{RQ1b: Disentangling Engine Identity from Sample Size}\label{subsec:rq1b}

To test whether PostgreSQL's higher accuracy in
Table~\ref{tab:main-results} reflects engine-specific structure or
simply its larger sample size, we apply the size-matched control
described in Section~\ref{subsec:sizecontrol-protocol}: PostgreSQL
subsampled to $n=200$ (class-proportional, 20 subsamples, each
evaluated under the same $10\times5$ repeated-CV protocol as the full
dataset, yielding 1000 fold-level estimates in total) is compared
against PostgreSQL at full size and against MySQL/MongoDB at their
native $n=200$.

\begin{table}[h]
\caption{Sample-Size Control: PostgreSQL Full vs.\ Size-Matched}\label{tab:size-control}
\begin{tabular}{@{}lrr@{}}
\toprule
Condition & Mean acc. & BCa 95\% CI \\
\midrule
PostgreSQL, full ($n=644$)         & 98.62\% & [98.40, 98.79] \\
PostgreSQL, matched ($n=200$)      & 96.69\% & [96.53, 96.86] \\
MySQL, native ($n=200$)            & 95.90\% & [95.30, 96.55] \\
MongoDB, native ($n=200$)          & 97.55\% & [96.90, 98.05] \\
\botrule
\end{tabular}
\end{table}

The size effect itself is present but more modest than a first look at
Table~\ref{tab:main-results} might suggest: once PostgreSQL is reduced
to the same sample size as MySQL and MongoDB, its accuracy drops by
$1.93$ percentage points (from $98.62\%$ to $96.69\%$), and its
confidence interval no longer overlaps with the full-size estimate
(Mann--Whitney $U=12556.0$, $p=6.7\times10^{-10}$). \textbf{This
confirms that part of the performance advantage attributed to
PostgreSQL in Section~\ref{subsec:rq1} is driven by dataset size, not
by engine identity --- though, as the comparison below shows, size is
not the whole story.}

The per-class breakdown (Table~\ref{tab:size-control-classes})
localises part of this effect in the PERFORMANCE class (the
sample-size-matched subsample also contains only $n=2$ AVAILABILITY
cases, too few to report a recall figure for that class here; see
Section~\ref{subsec:rq3} for the AVAILABILITY analysis): PERFORMANCE
recall drops from $97\%$ (full PostgreSQL) to $79.5\%$ (size-matched
PostgreSQL) once only 22 PERFORMANCE examples are available for
training instead of 70. This is a real effect of reduced training
data, but --- unlike what an earlier draft of this analysis reported
--- it does \emph{not} bring size-matched PostgreSQL down into the
same recall range as MySQL ($72\%$) or MongoDB ($58\%$); it remains
clearly above both. Sample size affects the PERFORMANCE class, but not
enough on its own to equalise it with the other two engines at the
same nominal training size. ECONOMIC remains at $100\%$ regardless of
sample size, consistent with its empirical separability discussed in
Section~\ref{subsec:economic-note}.

\begin{table}[h]
\caption{PERFORMANCE-Class Recall vs.\ Training Set Size}\label{tab:size-control-classes}
\footnotesize
\begin{tabular}{@{}p{3.2cm}rr@{}}
\toprule
Condition & $n_{\text{PERF.}}$ & Recall \\
\midrule
PostgreSQL, full ($n=644$)         & 70 & 97\% \\
PostgreSQL, size-matched ($n=200$) & 22 & 79.5\% \\
MySQL ($n=200$)                    & 18 & 72\% \\
MongoDB ($n=200$)                  & 12 & 58\% \\
\botrule
\end{tabular}
\end{table}

Notably, and revising what an earlier draft of this analysis
concluded, the size-matched PostgreSQL condition does not fall below
both other engines: it is significantly \emph{higher} than MySQL
native ($95.90\%$; Mann--Whitney $U=29797.0$, $p=0.017$) and
significantly \emph{lower} than MongoDB native ($97.55\%$;
$U=20607.0$, $p=0.028$), landing between the two rather than below
both. We interpret this pattern cautiously rather than as a settled
finding. Having confirmed (Limitation L1) that all three datasets are
produced by the same generator and the same random seed, a difference
in generation procedure is not a plausible explanation. The more
likely candidates are genuine finite-sample properties of how each
engine's particular scenario mix shapes the four-class decision
boundary at $n=200$, or sampling variance in the 20-subsample
protocol itself. We do not have grounds to attribute this ordering to
a property of the PostgreSQL, MySQL, or MongoDB \emph{engine} as
such, and flag it as an open question for future work
(Section~\ref{sec:conclusion}, FW4 below) rather than a settled
finding.

\begin{figure*}[t]
\centering
\includegraphics[width=0.85\linewidth]{figures/fig3_sample_size_control.pdf}
\caption{Sample-size control: mean cross-validated accuracy for
PostgreSQL at full size, PostgreSQL size-matched to $n=200$, and
MySQL/MongoDB at their native $n=200$. Error bars show BCa 95\%
confidence intervals (Table~\ref{tab:size-control}); the dashed line
marks full-size PostgreSQL accuracy for reference.}\label{fig:size-control}
\end{figure*}

\begin{figure*}[t]
\centering
\includegraphics[width=0.85\linewidth]{figures/fig4_performance_recall.pdf}
\caption{PERFORMANCE-class recall as a function of the number of
PERFORMANCE training examples available, across the four conditions
in Table~\ref{tab:size-control-classes}.}\label{fig:performance-recall}
\end{figure*}

\textbf{Implication for RQ1.} The evidence supports a revised answer
to RQ1, more measured than an earlier draft of this analysis
concluded: classifier performance is \emph{not} fully consistent
across engines at matched sample sizes, and training-set size is
\emph{part} of the explanation for the cross-engine differences
observed in Table~\ref{tab:main-results} --- narrowing PostgreSQL's
accuracy advantage over MySQL from $2.72$ points (native sample sizes)
to $0.79$ points once matched (PostgreSQL's own accuracy drops by
$1.93$ points once reduced to $n=200$), and clearly affecting
PERFORMANCE-class recall --- but size
alone does not fully equalise the three engines, since size-matched
PostgreSQL still lands between, rather than below, MySQL and MongoDB
native performance. Within the controlled synthetic design tested
here (Section~\ref{subsec:scope}), we conclude that both training-set
size and some remaining engine- or dataset-specific factor contribute
to the cross-engine differences, with size accounting for a
substantial but not exhaustive share. The practical implication is
correspondingly more modest: the data collection budget per class
likely matters for non-dominant modes such as PERFORMANCE, but it
should not be assumed to fully substitute for engine-specific
validation --- a conclusion we would not have reached without
re-running this analysis against the real underlying data, which we
flag as a caution about verifying computational results before
relying on them in a paper's central claims.

\subsection{RQ1c: Is RandomForest the Right Classifier?}\label{subsec:rq1c}

Every result so far uses RandomForest because it was the model
selected in \cite{iaha2026}, not because it was benchmarked against
alternatives for this specific task. Table~\ref{tab:classifier-comparison}
addresses this directly, comparing RandomForest against XGBoost, a
multinomial logistic regression, and an RBF-kernel SVM under the
identical protocol (Section~\ref{subsec:classifier-comparison-protocol}).

\begin{table}[h]
\caption{Classifier Comparison: Accuracy by Engine (Repeated $10\times5$-Fold CV)}\label{tab:classifier-comparison}
\begin{tabular}{@{}lrrr@{}}
\toprule
Classifier & MongoDB & MySQL & PostgreSQL \\
\midrule
RandomForest        & \textbf{97.55\%} & \textbf{95.90\%} & \textbf{98.62\%} \\
XGBoost              & 97.10\% & 95.35\% & 98.54\% \\
LogisticRegression   & 84.20\% & 84.35\% & 86.06\% \\
SVM (RBF)            & 83.65\% & 76.95\% & 77.16\% \\
\botrule
\end{tabular}
\end{table}

\begin{table}[h]
\caption{Classifier Comparison: Pooled 3-Class Macro-F1 by Engine}\label{tab:classifier-comparison-f1}
\begin{tabular}{@{}lrrr@{}}
\toprule
Classifier & MongoDB & MySQL & PostgreSQL \\
\midrule
RandomForest        & \textbf{93.15\%} & \textbf{93.83\%} & 98.93\% \\
XGBoost              & 90.99\% & 92.92\% & \textbf{98.97\%} \\
LogisticRegression   & 77.46\% & 84.08\% & 85.47\% \\
SVM (RBF)            & 75.20\% & 75.92\% & 77.86\% \\
\botrule
\end{tabular}
\end{table}

Two patterns are clear and, unlike several other results in this
paper, do not require heavy qualification. First, RandomForest is the
best or statistically indistinguishable from the best classifier on
every engine: it has the highest accuracy on all three engines, and
the highest macro-F1 on two of three (MongoDB, MySQL), with XGBoost
marginally ahead on PostgreSQL macro-F1 by $0.04$ percentage
points --- well within fold-to-fold noise. RandomForest was not,
therefore, an arbitrary or merely inherited choice: it is empirically
competitive with a stronger, more recent tree-ensemble alternative
(XGBoost) on exactly this task. Second, and more strikingly, both
tree-based methods (RandomForest, XGBoost) outperform both
linear/kernel methods (logistic regression, SVM) by a large and
consistent margin across all twelve tree-vs-linear/kernel pairings ---
from $11.0$ percentage points (XGBoost vs.\ logistic regression,
MySQL) to $21.5$ percentage points (RandomForest vs.\ SVM,
PostgreSQL) in accuracy. This is not a close contest. We read this as
consistent with, though not proof of, the underlying data-generating
process described in Limitation L1: labels are assigned by a
deterministic rule cascade over threshold conditions on the feature
values (Section~\ref{subsec:datasets}), which produces decision
boundaries that are naturally piecewise and axis-aligned --- exactly
the structure tree ensembles are suited to, and exactly what a global
linear decision boundary (logistic regression) or a smooth kernel
boundary (RBF-SVM) struggles to approximate. This does not mean
tree-based methods would necessarily dominate on real, non-synthetic
operational data with a different underlying structure, and we do not
claim otherwise; it does mean that, on the specific task and data
studied in this paper, the choice of RandomForest is empirically
justified rather than a default carried over unexamined from prior
work.

\subsection{RQ2: Cross-Engine Feature Importance Stability}\label{subsec:rq2}

\begin{table}[h]
\caption{Top-4 Feature Importances by Engine (RandomForest)}\label{tab:feature-importance}
\begin{tabular}{@{}lrrr@{}}
\toprule
Feature & MongoDB & MySQL & PostgreSQL \\
\midrule
\texttt{transactions\_per\_sec} & 0.227 & 0.279 & 0.238 \\
\texttt{cpu\_percent}           & 0.204 & 0.194 & 0.227 \\
\texttt{connections\_active}    & 0.150 & 0.149 & 0.171 \\
\texttt{replication\_lag\_mb}   & 0.122 & 0.135 & 0.158 \\
\botrule
\end{tabular}
\end{table}

\begin{figure*}[t]
\centering
\includegraphics[width=0.85\linewidth]{figures/fig2_feature_importance.pdf}
\caption{Top-4 RandomForest feature importances by engine
(Table~\ref{tab:feature-importance}). The same four features dominate
in the same relative order on all three engines.}\label{fig:feature-importance}
\end{figure*}

The same four features dominate the decision process across all three
engines, in the same relative order. Spearman rank correlation between
full feature-importance rankings (all ten features) is $\rho = 0.806$
(MongoDB vs.\ MySQL), $\rho = 0.915$ (MongoDB vs.\ PostgreSQL), and
$\rho = 0.891$ (MySQL vs.\ PostgreSQL). We emphasise what this
correlation does and
does not show: it indicates that the \emph{relative ranking} of which
features matter is stable across these three parametric variations of
a single generator (Section~\ref{subsec:scope}), which is a necessary
precondition for transfer learning to be worth attempting (there
would be little reason to transfer a model if the engines relied on
different signals entirely). It is not, by itself, evidence that a
classifier's learned decision boundary transfers to independently
instrumented engines --- that would require directly training on one
real engine and testing on another, which we identify as future work
(FW1, Section~\ref{sec:conclusion}) rather than claim here. We
therefore describe the decision signature as \emph{structurally
consistent} under this controlled design, and treat genuine
transferability as an open empirical question.

\subsection{A Note on ECONOMIC Separability}\label{subsec:economic-note}

Throughout this paper, the ECONOMIC class is perfectly separable
(recall and precision $=100\%$ on all three engines and at all sample
sizes tested), and our static heuristic baseline reproduces this
separation using a single threshold, $\texttt{cpu\_percent} < 30$
(Section~\ref{subsec:stats}). It is important to be precise about
\emph{why} this threshold works, since the underlying data-generation
process (consulted directly rather than inferred) reveals that
\texttt{cpu\_percent} is not itself the generating rule for ECONOMIC.
ECONOMIC samples are generated under a low-load
``off-peak'' scenario defined jointly by night-time hours and low
transaction throughput; CPU utilisation in that scenario happens to be
drawn from a low range ($10$--$30\%$) as a \emph{consequence} of the
same scenario, not as its defining criterion. The
$\texttt{cpu\_percent}<30$ threshold is therefore an effective
\emph{empirical proxy} for this scenario, not a description of the
causal mechanism that produced the label. We flag this distinction
explicitly because it affects how the perfect separability of
ECONOMIC should be read: it demonstrates that a single operational
metric can capture a multivariate, time-of-day-dependent operating
regime well enough to be perfectly diagnostic on this dataset, which
is a more interesting empirical finding than a merely definitional
threshold would be, but it should not be over-interpreted as evidence
that CPU utilisation alone \emph{causes} or \emph{defines} the
ECONOMIC mode in a deployed system.

\subsection{RQ3: The AVAILABILITY Failure Mode}\label{subsec:rq3}

Across all three engines, every AVAILABILITY misclassification goes to
CONSISTENCY, never to ECONOMIC or PERFORMANCE
(Table~\ref{tab:confusion-availability}). This directional pattern is
consistent across engines regardless of sample size and is the most
robust part of this result. The corresponding recall figure is $0\%$
in every engine, but we treat this number with caution rather than as
a precise estimate: with only $n=1$, $3$, and $5$ AVAILABILITY samples
respectively, recall is not statistically distinguishable from what a
single unlucky draw would produce, and we report it descriptively.

\begin{table}[h]
\caption{AVAILABILITY Misclassification Pattern by Engine}\label{tab:confusion-availability}
\begin{tabular}{@{}lrrrr@{}}
\toprule
Engine & $n_{\text{avail}}$ & Predicted CONS. & Recall & Feature overlap\footnotemark[2] \\
\midrule
MongoDB     & 1 & 1/1 & 0\% & n/a\footnotemark[3] \\
MySQL       & 3 & 3/3 & 0\% & 33.2\% \\
PostgreSQL  & 5 & 5/5 & 0\% & 38.9\% \\
\botrule
\end{tabular}
\footnotetext[2]{Mean \% of CONSISTENCY samples falling within the
observed AVAILABILITY feature range, averaged across the ten features.}
\footnotetext[3]{Not statistically interpretable with $n=1$.}
\end{table}

To test whether this failure is a correctable sample-size artefact or
a structural overlap, we applied SMOTE oversampling to the PostgreSQL
training folds (the only engine with enough AVAILABILITY samples,
$n=5$, to make oversampling meaningful). Because stratified 5-fold CV
with $n=5$ AVAILABILITY samples places exactly one such sample in each
held-out test fold, every training fold contains exactly four
AVAILABILITY samples; we therefore set SMOTE's neighbourhood-size
parameter to \texttt{k\_neighbors=3} (one below the number of
available minority
samples, the maximum value SMOTE permits) rather than its default of
5, which would not run with fewer than six minority samples per
training fold. SMOTE improved AVAILABILITY
recall from $0\%$ to only $20\%$ (F1 $=0.29$), with precision dropping
to $50\%$. This causal test was not repeated on MySQL or MongoDB, as
$n=1$ and $n=3$ AVAILABILITY samples are too few to support
oversampling ($n=3$ for MySQL would permit at most
\texttt{k\_neighbors=1} within a training fold containing only
2--3 such samples, which we judged too degenerate to be informative).
Combined with the descriptive feature-overlap evidence
that holds across all three engines ($33$--$39\%$ for MySQL and
PostgreSQL; not statistically interpretable for MongoDB at $n=1$), we
conclude that the failure is more consistent with a structural
overlap than with a purely volumetric shortage --- the current
ten-feature schema does not contain a feature that cleanly separates
AVAILABILITY from CONSISTENCY on the one engine where this could be
directly tested, and the same feature-space overlap is visible
wherever it can be measured. This is consistent with, and extends,
Limitation L3 of \cite{iaha2026}
(``CONSISTENCY $=89.4\%$ limits generalisation for ECONOMIC and
AVAILABILITY''), suggesting the same limitation is not specific to the
original PostgreSQL dataset; given the sample sizes involved, we frame
this as suggestive evidence pointing toward a structural explanation
rather than a definitive resolution of the question.

\subsection{RQ4: Live Recovery-Time Validation Across Engines}\label{subsec:rq4}

Table~\ref{tab:rto-results} reports total RTO (Section~\ref{subsec:live-validation})
for all three engines under the live failover protocol.

\begin{table}[h]
\caption{Total RTO by Engine (Live Kubernetes Failover Tests)}\label{tab:rto-results}
\begin{tabular}{@{}lrrrr@{}}
\toprule
Engine & $n$ & Mean total RTO & BCa 95\% CI & CV \\
\midrule
PostgreSQL & 30 & 37.26s & [33.20, 42.20] & 34.2\% \\
MySQL GR   & 14 & 105.29s & [104.96, 105.95] & 0.9\% \\
MongoDB RS & 20 & 43.82s & [41.06, 46.58] & 14.8\% \\
\botrule
\end{tabular}
\end{table}

All three pairwise Mann--Whitney comparisons are statistically
significant (PostgreSQL vs.\ MySQL: $U=0.0$, $p=1.3\times10^{-7}$,
rank-biserial $r=1.00$; PostgreSQL vs.\ MongoDB: $U=199.0$,
$p=0.047$, $r=0.34$; MySQL vs.\ MongoDB: $U=280.0$,
$p=1.1\times10^{-6}$, $|r|=1.00$), though given the orchestration
asymmetry stated in Section~\ref{subsec:live-validation}, we read
these as differences between \emph{deployments} rather than
differences intrinsic to the database engines themselves. The
rank-biserial correlation ($r=1-2U/(n_1 n_2)$) is large for two of
the three comparisons: PostgreSQL vs.\ MySQL and MySQL vs.\ MongoDB
both reach $|r|=1.00$, meaning every test in the faster engine's
sample recorded a lower total RTO than every test in the slower
engine's sample (complete stochastic dominance, not merely a
significant average difference). PostgreSQL vs.\ MongoDB is a more
moderate effect ($r=0.34$), consistent with their total RTO
distributions (Table~\ref{tab:rto-results}) overlapping more than either does
with MySQL's.

\textbf{MySQL Group Replication is slowest but by far the most
predictable.} At $105.29$s mean total RTO, MySQL GR is
$2.4$--$2.8\times$ slower than PostgreSQL or MongoDB in this setup,
but its coefficient of variation ($0.9\%$) is an order of magnitude
tighter than either alternative. This is consistent with MySQL Group
Replication's failure detection being dominated by fixed,
configuration-driven timeout parameters (e.g.\
\texttt{group\_replication\_member\_expel\_timeout} and related
settings) rather than by variable network- or consensus-driven
election dynamics, making its recovery behaviour highly predictable
even though slower in absolute terms.

\textbf{PostgreSQL shows a temporal ordering effect that we do not
fully explain.} Examining \texttt{promotion\_time\_s} in test order
reveals two clusters rather than random scatter: the first ten tests
average $34.02$s, while tests 11--30 average $13.69$s (with one late
outlier at test 23, $42.46$s). This pattern is consistent with the
CloudNativePG control loop, DNS propagation, or some other aspect of
cluster state becoming faster to converge after the first several
pod replacements in a run, but our test harness did not instrument
enough internal state to confirm the mechanism, and the 30 tests
should therefore not be treated as independent, identically
distributed draws from a single stationary RTO distribution --- a
point we return to in Threats to Validity
(Section~\ref{subsec:threats}).

\textbf{MongoDB's excluded fast cluster remains only partially
explained.} As described in Section~\ref{subsec:live-validation}, 10
of 30 MongoDB tests produced stabilisation times of $0.65$--$0.73$s,
tightly clustered and implausibly fast relative to the remaining 20
tests ($4.36$--$19.39$s). We exclude these 10 from
Table~\ref{tab:rto-results} rather than either silently dropping them
(as the original test script did for a different, unrelated reason)
or including them as genuine sub-second recoveries without comment.
We were not able to confirm the exact mechanism from the logs
available to us; the most plausible explanation, given that MongoDB
in this setup runs without a dedicated orchestration operator
(Section~\ref{subsec:live-validation}), is a stale replica-set status
read immediately following pod deletion, before Kubernetes had fully
evicted the terminating pod from the service's endpoint list. We
report this openly as an unresolved measurement question rather than
asserting a specific cause we cannot verify.

%% ============================================================
\section{Discussion}\label{sec:discussion}
%% ============================================================

\textbf{RQ1.} One result here is unambiguous and needs no hedging:
ML-driven mode selection significantly outperforms the static
CPU-threshold heuristic on every engine, by $18$--$21$ percentage
points, at $p<10^{-6}$ (Section~\ref{subsec:rq1}). Whatever else this
paper leaves open, that finding stands on its own. The more nuanced
question is \emph{why} accuracy varies across engines at all.
Classification accuracy is high across all three engines
but is not uniform, and our size-matched control
(Section~\ref{subsec:rq1b}) confirms \emph{empirically} --- rather
than merely suggesting --- that \emph{part} of this non-uniformity is
driven by training-set size, within the controlled synthetic design
tested here: subsampling PostgreSQL to MySQL/MongoDB's native size
narrows its accuracy advantage over MySQL from $2.72$ to $0.79$
points, drops PostgreSQL's own accuracy by $1.93$ points, and
noticeably reduces PERFORMANCE-class recall (from $97\%$ to
$79.5\%$). This is a useful, if more modest than initially expected,
signal for practitioners --- the data collection budget per class
(especially for non-dominant modes such as PERFORMANCE and
AVAILABILITY) plausibly matters, though evidently not enough on its
own to equalise engines at matched sample size --- and confirming
this on independently instrumented production engines remains to be
done (Section~\ref{subsec:scope}). Size does \emph{not} fully explain
the cross-engine differences: size-matched PostgreSQL lands between,
not below, MySQL and MongoDB accuracy (Section~\ref{subsec:rq1b}),
leaving a residual, dataset- or engine-specific effect that this
study is not positioned to fully explain.

\textbf{RQ2.} The strong correlation in feature importance across the
three engine variants ($\rho \geq 0.81$) suggests that
\texttt{transactions\_per\_sec}, \texttt{cpu\_percent},
\texttt{connections\_active}, and \texttt{replication\_lag\_mb}
constitute a largely engine-agnostic ``signature'' of replication-mode
state \emph{under this controlled design}. This is a necessary (though
not sufficient) condition for transfer learning across engines: it
justifies attempting a transfer-learning study --- pre-training on
PostgreSQL's larger dataset as a possible initialisation for MySQL or
MongoDB classifiers, particularly for the data-scarce PERFORMANCE and
AVAILABILITY classes --- but does not itself demonstrate that such
transfer would succeed on independently instrumented engines.

\textbf{RQ3.} The fact that AVAILABILITY is \emph{never} confused with
ECONOMIC or PERFORMANCE --- only ever with CONSISTENCY --- across all
three engines is itself informative: it suggests the current feature
schema captures a single dominant axis of variation (roughly,
load/CPU intensity) that separates ECONOMIC and PERFORMANCE well, but
does not capture whatever distinguishes an ``available-but-degraded''
state from a normal CONSISTENCY state. Candidate additional features
for future work include active replica/read-replica counts, health-check
failure rates, or failover-readiness indicators --- none of which are
present in the current ten-feature schema.

\subsection{Limitations}\label{subsec:limitations}

\textbf{L1 -- Synthetic data, single shared generator.} As
introduced in Section~\ref{subsec:datasets} and discussed at length in
Section~\ref{subsec:scope}, all three datasets are
synthetically generated by a single shared Python generator rather
than collected from live production traces, so this study should be
read as a controlled synthetic evaluation rather than an empirical
survey of real production systems (see also Threats to Validity,
External). We give the remaining mechanical detail here for
reproducibility. Each sample is drawn from one of five latent
operating scenarios (normal, peak, off-peak, crisis, degraded), each
with its own probability and its own parameter ranges for the ten
operational features; the four-class \texttt{mode} label is then
assigned by a deterministic rule cascade over the sampled feature
values (load thresholds, error rates, deadlock counts, and
time-of-day), not by a separate validation step applied to each
generated row. ``Expert-validated'' in this context and
in \cite{iaha2026} refers to the scenario probabilities, parameter
ranges, and labelling-rule cascade themselves having been specified
by domain expertise, not to an independent post-hoc review of each
individual sample. Two of the ten features
(\texttt{wal\_buffers\_full}, \texttt{trans\_rolled\_back}) are drawn
from engine-specific ranges within the shared generator (e.g.\
\texttt{wal\_buffers\_full} $\sim \mathcal{U}(0,50)$ for PostgreSQL
vs.\ $\mathcal{U}(0,20)$ for MongoDB), consistent with the
$2$--$2.5\times$ scale divergence we measured empirically between
engines on exactly these two features; all remaining features share
identical generating ranges across engines.

One consequence of this design is worth drawing out explicitly for
Section~\ref{subsec:rq1b}: because the same code and seed produce all
three datasets, a discrepancy \emph{between generators} cannot explain
the residual, unattributed ordering identified there (size-matched
PostgreSQL landing between MySQL and MongoDB at equal $n$, rather than
below both). This narrows, but does not resolve, the candidate
explanations for that residual; we discuss the remaining candidates
--- sampling variance and finite-sample structure specific to each
engine's scenario mix --- as an open question in FW4
(Section~\ref{sec:conclusion}) rather than a settled finding.

\textbf{L2 -- Extreme AVAILABILITY imbalance.} With $n=1$, $3$, and
$5$ AVAILABILITY samples respectively, no engine permits a
statistically reliable estimate of AVAILABILITY recall; we report it
descriptively rather than within aggregate metrics.

\textbf{L3 -- Non-independence of repeated-CV estimates.} Mann--Whitney
comparisons between engines use 50 fold-level accuracy estimates per
engine drawn from repeated resampling of the same underlying dataset;
these are not fully independent observations, so reported $p$-values
likely overstate true statistical confidence.

\textbf{L4 -- Live validation performed, but under an orchestration
asymmetry and with unresolved measurement questions.} We extend the
live Kubernetes validation of \cite{iaha2026} (Phase 1, CloudNativePG)
to MySQL and MongoDB (Section~\ref{subsec:rq4}), so this study is no
longer limited to offline classification performance. However, the
live comparison is confounded by an orchestration asymmetry
(PostgreSQL and MySQL run under production-grade operators; MongoDB
does not, Section~\ref{subsec:live-validation}), by a temporal
ordering effect in the PostgreSQL failover tests that we do not fully
explain, and by an unresolved measurement artefact affecting 10 of 30
MongoDB tests. The live RTO results should therefore be read as an
initial, imperfect validation rather than a definitive cross-engine
comparison.

\subsection{Threats to Validity}\label{subsec:threats}

\textbf{Internal.} The static heuristic baseline reuses the threshold
logic of IAHA's P1 layer rather than an independently tuned baseline
per engine, which may understate what a carefully engine-tuned
heuristic could achieve. Separately, the live PostgreSQL failover
tests (Section~\ref{subsec:rq4}) show a temporal ordering effect
(early tests substantially slower than later ones) that we cannot
fully attribute to a specific mechanism; the 30 tests should not be
treated as i.i.d.\ draws from a single stationary distribution, which
weakens the precision (though not necessarily the direction) of the
PostgreSQL RTO estimate.

\textbf{Construct.} Feature semantics (e.g.\ \texttt{wal\_buffers\_full}
as an oplog-buffer proxy for MongoDB) are functional analogues rather
than native engine metrics; this trades engine-specific fidelity for
cross-engine comparability, an explicit design choice discussed in
Section~\ref{sec:methodology}. Separately, the live RTO comparison
(Section~\ref{subsec:rq4}) is confounded by an orchestration
asymmetry: PostgreSQL and MySQL are managed by production-grade
Kubernetes operators (CloudNativePG, MySQL Operator), while MongoDB
runs as a bare StatefulSet with no dedicated operator. The reported
RTO differences therefore reflect \emph{deployment} differences at
least as much as they reflect properties of the database engines
themselves, and should not be read as a claim that MongoDB's
replication mechanism is inherently slower or less reliable than
PostgreSQL's or MySQL's.

\textbf{External.} As emphasised in Section~\ref{subsec:scope}, results
are obtained on synthetic datasets of differing size (200/200/644)
produced by a single shared generator, not on independently
instrumented production traces; generalisation to production traces,
to genuinely heterogeneous engine instrumentation, and to balanced
multi-engine datasets of equal size all remain to be validated (cf.\
L1, L2). This is the most significant threat to the external validity
of the paper's headline claims and should weigh heavily on how the
results are interpreted.

%% ============================================================
\section{Conclusion and Future Work}\label{sec:conclusion}
%% ============================================================

% [Confirmed by author, 19 Jul 2026: advisor (Prof. Assayad) has
%  reviewed and agrees with this Conclusion's framing.]
This study directly addresses FW4 from \cite{iaha2026} by evaluating,
in a controlled synthetic setting sharing a single generative process
across engines (Section~\ref{subsec:datasets}), whether ML-driven
replication mode selection behaves consistently across PostgreSQL,
MySQL, and MongoDB --- a necessary precondition for, though not a
demonstration of, cross-engine transfer (Section~\ref{subsec:scope}).
We find that classification accuracy is consistently high
and significantly outperforms static heuristics on all three engines
(McNemar $p<10^{-6}$), and that the underlying feature-importance
signature is highly stable across the three engine variants (Spearman
$\rho \geq 0.81$) --- evidence that supports, but does not yet
demonstrate, cross-engine transfer learning. A size-matched control
experiment shows that the cross-engine accuracy differences observed
at native dataset sizes are \emph{partly} --- but not fully --- driven
by training-set size rather than by intrinsic properties of the
database engine, within this controlled design: reducing PostgreSQL to
$n=200$ narrows its accuracy advantage over MySQL from $2.72$ to
$0.79$ points (PostgreSQL's own accuracy drops by $1.93$ points) and
noticeably affects PERFORMANCE-class recall, but size-matched
PostgreSQL still lands between, not below, MySQL and MongoDB, so a
residual difference between engines remains unexplained
(Section~\ref{subsec:rq1b}, FW4 below) and still requires validation
on independently instrumented engines before either the size effect or
the residual can be treated as an operational recommendation. At the same time, AVAILABILITY samples are consistently
misclassified as CONSISTENCY, never as ECONOMIC or PERFORMANCE, on all
three engines; because per-engine AVAILABILITY counts are very small
($n=1,3,5$), we treat this as a descriptive, hypothesis-generating
pattern rather than a settled recall estimate, and the diagnostic
causal test available on PostgreSQL points toward --- without fully
confirming --- a genuine gap in the current feature schema rather than
a simple data-volume artefact.

\textbf{FW1.} Explicit transfer-learning evaluation on independently
instrumented engines (pre-train on PostgreSQL, fine-tune on real
MySQL/MongoDB deployments), directly testing whether the
feature-importance stability found here under a shared generator
translates into sample-efficiency gains for the data-scarce engines
in practice.

\textbf{FW2.} Extend the feature schema with availability-specific
signals (replica health-check failure rate, active read-replica count)
to test whether the structural AVAILABILITY/CONSISTENCY overlap
identified in Section~\ref{subsec:rq3} can be resolved.

\textbf{FW3.} Resolve the two open measurement questions from the live
validation in Section~\ref{subsec:rq4}: confirm the mechanism behind
the excluded fast MongoDB cluster (most plausibly a stale status read
during pod eviction), and re-run PostgreSQL failover tests with
sufficient repetitions and instrumentation to determine why early
tests are consistently slower than later ones. Repeat the live
validation with MongoDB under a dedicated Kubernetes operator to
remove the orchestration asymmetry noted in
Section~\ref{subsec:live-validation}, so that engine differences can
be separated from deployment differences. Extend the preliminary
RPO/data-durability piloting we have begun on MongoDB to PostgreSQL
and MySQL under the same protocol.

\textbf{FW4.} Investigate the residual, unattributed ordering between
size-matched PostgreSQL, MySQL, and MongoDB identified in
Section~\ref{subsec:rq1b} (size-matched PostgreSQL lands significantly
above MySQL but significantly below MongoDB at the same $n=200$).
Having confirmed (Limitation L1) that all three datasets share a
single generator and seed, a generator discrepancy is not a candidate
explanation; remaining hypotheses are sampling variance not fully
captured by our 20-subsample protocol, or genuine finite-sample
structure in how the four-class decision boundary behaves under each
engine's particular scenario mix at $n=200$ specifically.

%% ============================================================
\backmatter

\bmhead{Acknowledgments}
Not applicable.

\section*{Statements and Declarations}

% [Confirmed by author, 19 Jul 2026: no funding, no competing
%  interests, no ethics/consent issues -- statements below stand as
%  accurate.]
\begin{itemize}
\item \textbf{Funding:} The authors did not receive support from any
  organization for the submitted work.
\item \textbf{Competing Interests:} The authors have no relevant
  financial or non-financial interests to disclose.
\item \textbf{Ethics approval:} Not applicable.
\item \textbf{Consent to participate:} Not applicable.
\item \textbf{Consent for publication:} Not applicable.
\item \textbf{Availability of data and materials:} The datasets
  generated and analysed in this study, together with the data
  generation script, are publicly available in the GitHub repository
  \url{https://github.com/imanejerrari-etu-bit/cross-db-replication-mode-selection}
  under a CC-BY-4.0 license.
  % [Author-confirmed, 19 Jul 2026: all analysis scripts
  %  (train_classifiers.py, stats_tests.py, sample_size_control.py,
  %  smote_availability.py, analyze_rto_total.py, common.py) have
  %  been pushed to the repository. Self-reported by the author; not
  %  independently re-verified in this session (no web access
  %  available at the time of this edit). Spot-check the live repo
  %  once before final submission to be certain.]
\item \textbf{Code availability:} The analysis code used to produce
  the results in this paper, together with the dataset generation
  script, is available in the same repository
  (\url{https://github.com/imanejerrari-etu-bit/cross-db-replication-mode-selection})
  under an MIT license.
\item \textbf{Authors' contributions:} All authors contributed to the
  study conception and design. Material preparation, data collection,
  and analysis were performed by I.~Jerrari. The first draft of the
  manuscript was written by I.~Jerrari and all authors commented on
  previous versions of the manuscript. All authors read and approved
  the final manuscript.
  % [Confirmed by author, 19 Jul 2026: the sole co-author, Prof.
  %  Assayad, is also the thesis supervisor and has already agreed to
  %  this paper's framing (see Conclusion sign-off above); this
  %  contribution statement stands as accurate.]
\item \textbf{Use of AI tools:} A Large Language Model (LLM) tool
  assisted with manuscript preparation, including copy editing and
  reviewing the scope and framing of certain claims. All experiments,
  analyses, results, and interpretations are the authors' own, who
  take full responsibility for the manuscript's content. No
  generative AI tool is listed as an author.
\end{itemize}

% ---- Bibliography (numbered) ----
\begin{thebibliography}{20}

\bibitem{iaha2026}
Jerrari, I., Assayad, I.: IAHA: An intelligent adaptive high
availability framework for containerised PostgreSQL databases using
hybrid machine learning. Journal of Intelligent \& Fuzzy Systems,
under review (2026)

\bibitem{patroni}
Patroni Project: Patroni: A template for PostgreSQL HA.
\url{https://github.com/patroni/patroni}. Accessed Apr 2025

\bibitem{repmgr}
repmgr Project: repmgr: PostgreSQL replication manager.
\url{https://repmgr.org/}. Accessed Apr 2025

\bibitem{stolon}
Sorint.lab: Stolon: PostgreSQL cloud-native HA.
\url{https://github.com/sorintlab/stolon}. Accessed Apr 2025

\bibitem{cloudnativepg}
EnterpriseDB/CNCF: CloudNativePG: Cloud-native PostgreSQL operator.
\url{https://cloudnative-pg.io/}. Accessed Apr 2025

\bibitem{kraska2019}
Kraska, T., et al.: SageDB: A learned database system. In: Proc. CIDR
(2019)

\bibitem{marcus2019}
Marcus, R., Negi, P., Mao, H., Zhang, C., Alizadeh, A., Kraska, T.,
Papaemmanouil, O., Tatbul, N.: Neo: A learned query optimizer. Proc.
VLDB Endow. \textbf{12}(11), 1705--1718 (2019)

\bibitem{marcus2021}
Marcus, R., et al.: Bao: Making learned query optimization practical.
In: Proc. ACM SIGMOD, pp.\ 1275--1288 (2021)

\bibitem{vanaken2017}
Van Aken, D., Pavlo, A., Gordon, G.J., Zhang, B.: Automatic database
management system tuning through large-scale machine learning. In:
Proc. ACM SIGMOD, pp.\ 1009--1024 (2017).
\url{https://doi.org/10.1145/3035918.3064029}

\bibitem{zhang2019}
Zhang, J., et al.: An end-to-end automatic cloud database tuning
system using deep reinforcement learning. In: Proc. ACM SIGMOD, pp.\
415--432 (2019)

\bibitem{autopilot2020}
Rzadca, K., Findeisen, P., Swiderski, J., Zych, P., Broniek, P.,
Kusmierek, J., Nowak, P., Strack, B., Witusowski, P., Hand, S.,
Wilkes, J.: Autopilot: Workload autoscaling at Google. In: Proc.\ 15th
Eur.\ Conf.\ Comput.\ Syst.\ (EuroSys '20), Heraklion, Greece, Article
no.\ 16 (2020). \url{https://doi.org/10.1145/3342195.3387524}

\bibitem{firm2020}
Qiu, H., et al.: FIRM: An intelligent fine-grained resource
management framework for SLO-oriented microservices. In: Proc.\ OSDI,
pp.\ 805--825 (2020)

\bibitem{gunning2019}
Gunning, D., et al.: XAI---Explainable artificial intelligence. Sci.
Robot. \textbf{4}(37), eaay7120 (2019)

\bibitem{lundberg2017}
Lundberg, S.M., Lee, S.-I.: A unified approach to interpreting model
predictions. In: Proc.\ NeurIPS, pp.\ 4765--4774 (2017)

\bibitem{ribeiro2016}
Ribeiro, M.T., Singh, S., Guestrin, C.: ``Why should I trust you?'':
Explaining the predictions of any classifier. In: Proc.\ ACM KDD, pp.\
1135--1144 (2016)

\bibitem{wachter2017}
Wachter, S., Mittelstadt, B., Russell, C.: Counterfactual explanations
without opening the black box: Automated decisions and the GDPR.
Harvard J. Law Technol. \textbf{31}(2), 841--887 (2017)

\bibitem{guidotti2018}
Guidotti, R., et al.: A survey of methods for explaining black box
models. ACM Comput. Surv. \textbf{51}(5), 93:1--93:42 (2018).
\url{https://doi.org/10.1145/3236009}

\bibitem{holzinger2019}
Holzinger, A., Langs, G., Denk, H., Zatloukal, K., M\"uller, H.:
Causability and explainability of artificial intelligence in
medicine. WIREs Data Min. Knowl. Discov. \textbf{9}(4), e1312 (2019).
\url{https://doi.org/10.1002/widm.1312}

\bibitem{atakishiyev2021}
Atakishiyev, S., Salameh, M., Yao, H., Goebel, R.: Explainable
artificial intelligence for autonomous driving: a comprehensive
overview and field guide for future research directions. IEEE
Access \textbf{12}, 101603--101625 (2024).
\url{https://doi.org/10.1109/ACCESS.2024.3431437}

\bibitem{bussmann2020}
Bussmann, N., et al.: Explainable machine learning in credit risk
management. Comput. Econ. \textbf{57}, 203--216 (2021)

\bibitem{kanellis2022}
Kanellis, K., Ding, C., Kroth, B., M\"uller, A., Curino, C.,
Venkataraman, S.: LlamaTune: Sample-efficient DBMS configuration
tuning. Proc. VLDB Endow. \textbf{15}(11), 2953--2965 (2022).
\url{https://doi.org/10.14778/3551793.3551844}

\bibitem{lao2024}
Lao, J., Wang, Y., Li, Y., Wang, J., Zhang, Y., Cheng, Z., Chen, W.,
Tang, M., Wang, J.: GPTuner: A manual-reading database tuning system
via GPT-guided Bayesian optimization. Proc. VLDB Endow.
\textbf{17}(8), 1939--1952 (2024).
\url{https://doi.org/10.14778/3659437.3659449}

\bibitem{wang2021udo}
Wang, J., Trummer, I., Basu, D.: UDO: Universal database optimization
using reinforcement learning. Proc. VLDB Endow. \textbf{14}(13),
3402--3414 (2021)

\end{thebibliography}

\end{document}
